\par\noindent\textbf{(1)} 在 \(x=\pm a\) 处波函数连续，导数跳跃为 \(\psi'(a^+)-\psi'(a^-)=-2mg\psi(a)/\hbar^2\)，\(-a\) 处相同。势关于原点偶对称，可将束缚态分别取偶宇称与奇宇称；外区解按 \(e^{-\kappa|x|}\) 衰减。
\par\noindent\textbf{(2)} 偶态在 \(|x|<a\) 取 \(\psi\propto\cosh(\kappa x)\)，外区衰减。于 \(x=a\) 连接后 \(-\kappa-\kappa\tanh(\kappa a)=-2mg/\hbar^2\)，整理为 \(\kappa=(mg/\hbar^2)[1+e^{-2\kappa a}]\)。对任意 \(g>0\) 至少有一个正偶态根。
\par\noindent\textbf{(3)} 奇态在 \(|x|<a\) 取 \(\psi\propto\sinh(\kappa x)\)，连接给 \(-\kappa-\kappa\coth(\kappa a)=-2mg/\hbar^2\)，即 \(\kappa=(mg/\hbar^2)[1-e^{-2\kappa a}]\)。当 \(\kappa\to0\) 时右侧斜率为 \(2amg/\hbar^2\)，因此奇束缚态仅在 \(2amg/\hbar^2>1\) 时存在。
\par\noindent\textbf{(4)} 令 \(\kappa_0=mg/\hbar^2\)、\(\kappa_0a\gg1\)，偶奇根近似 \(\kappa_\pm\approx\kappa_0[1\pm e^{-2\kappa_0a}]\)。能量 \(E_\pm=-\hbar^2\kappa_\pm^2/(2m)\)，劈裂模约 \(|E_+-E_-|\approx(2\hbar^2\kappa_0^2/m)e^{-2\kappa_0a}\)，体现量子隧穿的指数衰减。
